\section{Reconstructing the original hypersurface}
\label{emb-section}

The static isometry in \cref{hw-static-rigidity} concerns the quotient
metric.  We now reconstruct the original data, including their original
smooth structure at the boundary and their future second fundamental
form.  Put
\begin{equation}
 M=m_{\AH}=\frac{r_h+r_h^3}{2},\qquad
 F(r)=1+r^2-\frac{2M}{r},\qquad
 q=\frac{dr^2}{F(r)}+r^2\sigma .
 \label{emb-model}
\end{equation}
The theorem provides an isometry of static systems
\[
 \Psi:(\operatorname{int}\Omega,h,\lambda)
       \longrightarrow ((r_h,\infty)\times S^2,q,\sqrt F)
\]
which extends to their complete Riemannian boundaries.  We use
$(r,\omega)$ for its coordinates.  Notice that
\[
 F'(r)=2r+\frac{2M}{r^2}>0,\qquad F'(r_h)=2\kappa .
\]
Thus $F$ has a smooth inverse near its zero, and
$r-r_h=N/(2\kappa)+O(N^2)$, where $N=\lambda^2$.

\subsection{The interior graph and its future normal}

\begin{lemma}[The original interior data]
\label{emb-interior}
There is a smooth function $H$ on $\operatorname{int}\Omega$ for
which
\begin{equation}
 I(x)=\bigl(T=H(x),\,r(x),\,\omega(x)\bigr)
 \label{emb-graph}
\end{equation}
is a spacelike embedding into the exterior with metric
\[
 \mathbf g_M=-F(r)\,dT^2+\frac{dr^2}{F(r)}+r^2\sigma .
\]
Its induced metric is exactly $g$, and its second fundamental form
with the future normal is exactly $K$.
\end{lemma}

\begin{proof}
By \cref{geo-static-quotient}, the closed one-form
$\mathcal A=N^{-1}X^\flat$ has a global primitive with
$dH=-\mathcal A$.  Set $T=t+H(x)$ on
$\mathbb R\times\operatorname{int}\Omega$.  Completion of the square
gives the exact identity
\begin{equation}
 h-N(dT)^2
 =-u^2dt^2+g_{ij}(dx^i+X^i dt)(dx^j+X^j dt).
 \label{emb-stationary-identity}
\end{equation}
After applying $\Psi$ on the spatial factor, the left side is
$\mathbf g_M$.  The slice $t=0$ is precisely \eqref{emb-graph}.
In particular,
\[
 I^*\mathbf g_M=h-N\,dH\otimes dH
              =h-N^{-1}X^\flat\otimes X^\flat=g.
\]
The projection of the graph to the static base is the diffeomorphism
$\Psi$, so the graph is injective and is an embedding in the interior.

Choose the time orientation for which $\partial_T=\partial_t$ is
future timelike in this exterior.  The unit normal to $t=0$ is
\[
 n=u^{-1}(\partial_t-X),
 \qquad
 \mathbf g_M(n,\partial_t)=-u<0,
\]
so it is future directed.  With the convention
$K_{ij}=\mathbf g_M(\overline\nabla_i n,\partial_j)$, the
ADM identity is
$\partial_tg=2uK+\mathcal L_Xg$.
The metric on the right of \eqref{emb-stationary-identity} is
stationary, and \cref{eq-causal-killing} gives
$\mathcal L_Xg=-2uK$ for the original tensor.  Hence the induced
future second form agrees with that original tensor.
\end{proof}

\subsection{Regular horizon coordinates}

Fix any antiderivative $r_*$ of $F^{-1}$ on $(r_h,\infty)$.
The simple root gives
\begin{equation}
 r_*=\frac1{2\kappa}\log(r-r_h)+c_*(r),
 \qquad
 e^{\kappa r_*}=\sqrt{F(r)}\,Q(F(r)),
 \label{emb-tortoise}
\end{equation}
where $c_*$ is smooth near $r_h$ and $Q$ is a positive smooth
function near $0$.  Define exterior Kruskal coordinates by
\begin{equation}
 U=-e^{-\kappa(T-r_*)},\qquad
 V=e^{\kappa(T+r_*)}.
 \label{emb-kruskal}
\end{equation}
Then $U<0<V$ in this exterior and
\begin{equation}
 -UV=FQ(F)^2,\qquad
 \mathbf g_M=-\frac1{\kappa^2Q(F)^2}\,dU\,dV+r^2\sigma .
 \label{emb-kruskal-metric}
\end{equation}
Here $dU\,dV$ denotes the symmetric product.  The derivative of
$FQ(F)^2$ with respect to $r$ at $r_h$ is
$2\kappa Q(0)^2>0$; hence $r$ and the coefficients in
\eqref{emb-kruskal-metric} are smooth functions of $-UV$ near
the horizon.  This gives regular coordinates in the smooth maximal
extension.  Its future horizon adjoining our exterior is
$U=0$, $V>0$, with the bifurcation sphere at $U=V=0$.

\begin{lemma}[Smoothness in the original boundary structure]
\label{emb-horizon}
The graph \eqref{emb-graph} extends smoothly to the original closed
manifold $\Omega$.  If $u|_S>0$, its boundary is a full smooth
cross-section of $U=0$, $V>0$.  If $u|_S=0$, its boundary is
the bifurcation sphere.  In both cases the boundary map is a
diffeomorphism onto that section.
\end{lemma}

\begin{proof}
First suppose $u|_S>0$.  By \cref{geo-boundary-collars},
$N$ is a smooth defining function in the original collar and
\begin{equation}
 \mathcal A=\frac1{2\kappa}\,d\log N+\beta,\qquad
 H=-\frac1{2\kappa}\log N+B ,
 \label{emb-positive-primitive}
\end{equation}
where $\beta$ is smooth and closed up to $S$ and $B$ is smooth
there.  To justify the last assertion directly, the differential of
the displayed remainder in $H$ is $-\beta$.  Integrating its bounded
normal derivative gives a boundary value; tangential derivatives
then extend by the smooth differential equation.  This works to
every order and on overlapping collar patches.

The angular coordinates of $\Psi$ are also smooth in the original
collar.  Here is the point that is needed because the regular base
uses $\lambda=\sqrt N$.  Formula \eqref{st-even-collar} doubles $h$
across $\lambda=0$ by the isometry
$(\lambda,y)\mapsto(-\lambda,y)$.  The angular coordinate $\omega$
of the model is its normal footpoint on the horizon: in the model,
curves of constant $\omega$ are the normal geodesics from the
boundary.  A boundary-preserving isometry has the same description
on the base.  Reflection takes each such geodesic to its continuation
with the same footpoint.  Therefore $\omega$ is a smooth even
function of $\lambda$ in doubled collar coordinates.  A smooth
even function of one real variable, with smooth parameters, is a
smooth function of its square on the half-line.  This follows, to
each finite order, from Taylor expansion with its smooth remainder.
Thus $\omega=\omega(N,y)$ is smooth in the original collar.
Likewise $r=F^{-1}(N)$ is smooth.

Substituting \eqref{emb-positive-primitive} into
\eqref{emb-tortoise}--\eqref{emb-kruskal} gives on the graph
\begin{equation}
 U=-Q(N)\,N e^{-\kappa B},\qquad
 V=Q(N)e^{\kappa B}.
 \label{emb-positive-kruskal}
\end{equation}
These are smooth original collar functions.  On $S$ they satisfy
$U=0$ and $V=Q(0)e^{\kappa B}>0$.  The restriction
$\omega:S\to S^2$ is a diffeomorphism because $\Psi$ extends
as a boundary isometry.  Hence the image is the graph of a smooth
positive function $V(\omega)$ over the whole horizon sphere, meeting
each generator once.  The horizon first form is $r_h^2\sigma$,
so this section is spacelike.

Now suppose $u|_S=0$.  In the original normal coordinate $s\geq0$,
\cref{geo-boundary-collars} gives smooth fields
\[
 u=sa,\quad X=s^2Y,\quad
 N=s^2D,\quad D=a^2-s^2|Y|_g^2,\quad D|_S=\kappa^2,
 \qquad \mathcal A=D^{-1}Y^\flat .
\]
The base smooth structure is the original one, and $H$ extends
smoothly by the same primitive argument.  On this one-sided collar
$\lambda=s\sqrt D$ is smooth, with positive normal derivative.
Both the angular boundary map and $r=F^{-1}(s^2D)$ are smooth.
Equations \eqref{emb-kruskal} become
\begin{equation}
 U=-Q(N)\,\lambda e^{-\kappa H},\qquad
 V=Q(N)\,\lambda e^{\kappa H}.
 \label{emb-zero-kruskal}
\end{equation}
They extend smoothly with $U=V=0$ on $S$ and with nonzero
normal derivatives of opposite signs.  The angular map again gives
a diffeomorphism of $S$ onto the entire bifurcation sphere.

In either case the extended map has induced metric $g$ at $S$
by continuity of the identity in \cref{emb-interior}, since all
ambient and map coefficients just displayed are smooth.  For any
nonzero tangent vector $v$ of $\Omega$, including at its boundary,
\[
 \mathbf g_M(dI(v),dI(v))=g(v,v)>0.
\]
Thus $dI$ is injective and its image is spacelike up to the boundary.
This establishes smoothness and immersion in the original, rather
than merely the completed-base, structure.
\end{proof}

\begin{proposition}[Properness, the future second form, and infinity]
\label{emb-global}
The extended map $I$ is a smooth proper spacelike embedding of the
original $(\Omega,g,K)$ into the smooth maximal Schwarzschild--AdS
extension of parameter $M=m_{\AH}$.  Its interior lies in one
exterior, its boundary is the full future horizon section described
in \cref{emb-horizon}, and its end reaches that exterior's conformal
infinity.
\end{proposition}

\begin{proof}
The interior map is injective by its base projection, the boundary
map is injective by \cref{emb-horizon}, and their images are
disjoint since $r>r_h$ in the interior and $r=r_h$ on $S$.
Let $C$ be a compact subset of the maximal spacetime.  Its radius
function is bounded above, say by $R$.  The preimage $I^{-1}(C)$
is closed and lies in
\[
 \Psi^{-1}\bigl([r_h,R]\times S^2\bigr),
\]
which is compact in the completed base.  That completion is
homeomorphic to the original closed exterior, so it is also compact
in $\Omega$.  Hence $I$ is proper.  A proper injective immersion is
an embedding, including for manifolds with boundary.

The future unit normal extends smoothly to $S$.  One can construct
it in any of the regular charts by choosing a smooth future timelike
ambient vector, projecting it onto the normal line of the smooth
spacelike tangent bundle, and normalizing.  These local choices
give the unique future normal and therefore agree on overlaps.
In the interior it is the normal used in \cref{emb-interior}.
Its second form is smooth up to $S$ and agrees with the original
$K$ in the interior; continuity proves equality at $S$ as well.

For the end statement, temporarily denote the original end radius
by $r_o$.  The original estimates give
\[
 |\mathcal A|_g
 =\frac{|X|_g}{N}=O(r_o^{-1-\tau_1}),\qquad
 \tau_1>3/2.
\]
Consequently, in the original radial and angular coordinates,
\[
 \mathcal A_{r_o}=O(r_o^{-2-\tau_1}),\qquad
 \mathcal A_A=O(r_o^{-\tau_1}).
\]
Radial integration of $dH=-\mathcal A$ gives a uniform limit of
$H$ as $r_o\to\infty$.  Its angular oscillation tends to zero by
the second bound, so this limit is a single constant $T_\infty$.
Moreover $\lambda/\sqrt{1+r_o^2}\to1$ and $\lambda=\sqrt{F(r)}$,
which imply $r/r_o\to1$.  The unique end therefore has
$r\to\infty$ and $T\to T_\infty$.  The base isometry covers a
full model end, so all angles occur.  In the standard conformal
completion with defining function $1/r$, the graph thus reaches
the full section $T=T_\infty$ of this exterior's conformal infinity.
All these conclusions follow from the reconstructed graph and its
specific decay estimates.
\end{proof}

\subsection{The converse and explicit equality examples}

\begin{proposition}[Converse on the stated horizon subclass]
\label{emb-converse}
Suppose the original data satisfy the connected horizon hypotheses
of \cref{def:setup-horizon} and have the embedding required in
\cref{thm:main-rigidity}, with its stipulated actual metric mass
$M=m_{\AH}$.  Then equality holds in the Penrose Inequality.
\end{proposition}

\begin{proof}
In a regular horizon chart the first form of $U=0$ has null
generator $\partial_V$ in its kernel and angular part
$r_h^2\sigma$.  Every full smooth spacelike cross-section therefore
has area $4\pi r_h^2$, including the bifurcation section.  The
embedding and the assumed outer area-minimization give
\[
 A_{\min}=\Area_g(S)=4\pi r_h^2.
\]
The equation $F_M(r_h)=0$ and the actual mass match now yield
\[
 m_{\AH}=M=\frac{r_h+r_h^3}{2}
 =\sqrt{\frac{A_{\min}}{16\pi}}
  \left(1+\frac{A_{\min}}{4\pi}\right).
\]
This uses the mass match as specified in the statement.  No
invariance under an arbitrary change of asymptotic time slice is
needed for this implication.
\end{proof}

Propositions~\ref{emb-global} and \ref{emb-converse} prove
Theorem~\ref{thm:main-rigidity}. The following examples also verify
sharpness without a maximality assumption.

\begin{proposition}[Mass-matched static and nonstatic examples]
\label{emb-examples}
For every $M>0$, the static exterior with its bifurcation boundary
attains equality.  It also admits smooth compactly supported time
graph perturbations with $K\not\equiv0$ which satisfy all the
initial-data and connected horizon hypotheses and still attain
equality with actual mass $m_{\AH}=M$.
\end{proposition}

\begin{proof}
Fix a closed annulus
$r_h<a<b<\infty$ and a smooth function $\varphi$ supported in
its interior.  On the graph $T=\varphi(r,\omega)$ the induced
metric is
\begin{equation}
 q_\varphi=q-F\,d\varphi\otimes d\varphi .
 \label{emb-example-metric}
\end{equation}
For sufficiently small $C^2$ norm of $\varphi$ it is positive
and uniformly comparable with $q$, so its completion is a smooth
complete exterior with the same boundary.  In proper distance from
the horizon, $q$ is smooth and has a totally geodesic boundary;
$\varphi=0$ on that collar.  The graph lies in the vacuum metric
$\Ric_{\mathbf g_M}=-3\mathbf g_M$, so its data satisfy the vacuum
constraints and hence the dominant energy condition.  Near the
boundary and near infinity its future second form vanishes and
its metric is exactly $q$.

Every coordinate sphere with $r>r_h$ has
\[
 H_q=\frac{2\sqrt F}{r}>0.
\]
Its future expansion for the perturbed graph remains positive on
$[a,b]$ by $C^2$ smallness, uniformly on this fixed annulus, and
is unchanged outside it.  If a weakly future trapped enclosing
surface had a component strictly outside $S$, take its maximum
radius.  There it touches a coordinate sphere from the smaller
radius side, with the same outward normal toward the end.  The
mean curvature comparison gives $H_{\rm surface}\geq
H_{\rm sphere}$, while the tangential planes and hence the
tangential traces of $K$ agree.  Its future expansion would be
positive at that point, a contradiction.  Thus the horizon is
outermost for the required future condition.

Let $\pi:[r_h,\infty)\times S^2\to S^2$ be angular projection.
On the perturbation annulus the quadratic form
$q-r_h^2\pi^*\sigma$ is positive definite with a uniform lower
bound.  Taking $\varphi$ smaller if necessary gives everywhere
\[
 q_\varphi\geq r_h^2\pi^*\sigma.
\]
For any admissible full enclosing cut $\Gamma=\partial_\Omega D$,
including cuts with portions on $S$, apply Stokes' Theorem to a
large truncation of $D$ and the closed form
$\pi^*dA_\sigma$.  With the orientation toward the chosen end,
the sum of the degrees of all its boundary components is one:
\[
 \int_\Gamma \pi^*dA_\sigma=4\pi .
\]
This uses the full intrinsic boundary, so when $D=\Omega$ the
integral is over $S$ itself.  The pointwise quadratic-form bound
gives
\[
 \Area_{q_\varphi}(\Gamma)
 \geq r_h^2\int_\Gamma|\pi^*dA_\sigma|
 \geq4\pi r_h^2.
\]
Since $S$ has that area, it is outer area-minimizing.

We compute the actual metric flux.  The metric equals $q$ near
infinity, so put $r=\sinh\eta$ and $V_0=\cosh\eta$.  Relative
to $b=d\eta^2+r^2\sigma$ its only nonzero error component is
\begin{equation}
 q-b=f\,d\eta^2,\qquad
 f=\frac{2M}{rF(r)}
   =2Mr^{-3}+O(r^{-5}).
 \label{emb-example-error}
\end{equation}
Writing $\nu_b=\partial_\eta$, direct background differentiation
gives
\[
 (\operatorname{div}_b(q-b)-d\operatorname{tr}_b(q-b))(\nu_b)
       =\frac{2V_0}{r}f.
\]
For any of the four static potentials $V$, the radial component
of
$\operatorname{tr}_b(q-b)\,dV-(q-b)(\nabla_bV,\cdot)$
vanishes.  Thus the time-component flux in
\cref{def:setup-mass} is
\begin{equation}
 p_0=\lim_{r\to\infty}
       \frac{1}{16\pi}\,4\pi r^2\,
       \frac{2V_0^2}{r}f
     =\lim_{r\to\infty}M\,\frac{V_0^2}{F(r)}
     =M.
 \label{emb-example-mass}
\end{equation}
For $V_i=r\omega_i$ the remaining integrand is a radial multiple
of $\omega_i$, whose spherical integral vanishes.  Hence
$(p_0,p_1,p_2,p_3)=(M,0,0,0)$ and $m_{\AH}=M$.
Equation \eqref{emb-example-error} gives the required differentiated
decay; all errors involving $K$, and all scalar-curvature defects
from $R=-6$, are compactly supported.  The regularity and weighted
integrability assumptions follow.

The graph is smooth at the bifurcation boundary and proper by the
argument of \cref{emb-global}, or directly because it is unchanged
there and outside a compact annulus.  Thus all the conditions in
\cref{emb-converse} hold.  Choosing $\varphi$ with a nonzero Hessian
at a critical point gives $K\neq0$ there: at such a point the graph
second form is $\sqrt F\,\Hess_q\varphi$.  These furnish
non-time-symmetric equality examples as well as the static one.
\end{proof}
